\section{評価結果}\label{chapter:eval}

\subsection{評価する問いと証拠の種類}\label{sec:eval:sample}

評価では四つの問いを分ける。\emph{適法性}は、不正な制度変更を拒否し、
正しい変更を不可分な一単位として確定できるかを問う。\emph{活動性}は、
提案、攻撃、裁定、遷移、修復、再試験までの処理が進むかを問う。
\emph{検出性能}は、既知の故障をどの程度見つけ、正常例をどの程度誤って
拒否するかを問う。\emph{有効性}は、統治された共進化が研究品質、評価
の校正、費用を改善するかを問う。試験は主に適法性を、固定故障例は特定
経路の活動性を調べる。検出性能と有効性には、正解付き事例、独立試行、
対照条件が必要である。

\begin{table}[t]
  \centering
  \small
  \begin{tabularx}{\columnwidth}{@{}lXX@{}}
    \toprule
    \textbf{証拠} & \textbf{支持できる主張} & \textbf{支持できない主張} \\
    \midrule
    1,124 件の関連試験 & 実装した規則が検査例どおり働くこと &
      未試験入力での安全性、研究品質 \\
    五故障・一正常例 & 指定した五経路の発火と正常例の通過 &
      誤検出率、見逃し率、未知故障への一般化 \\
    探索的な動作記録 & 一部の長い処理経路を実行できたこと &
      再現性、分布全体での活動性、優越性 \\
    未実施の対照実験 & なし &
      品質、校正、費用、構成要素ごとの因果効果 \\
    \bottomrule
  \end{tabularx}
  \caption{現時点の証拠から支持できる主張と、まだ支持できない主張。}
  \label{tab:eval-status}
\end{table}

\subsection{一連の動作確認}\label{sec:eval:construction}

成果物の説明に記録した対象選択条件では、RQGM と最終原稿検証に直接
関係する 1,124 件を収集し、この作業版で全件が合格した。全体の選択条件
では 4,805 件を収集し、4,787 件が合格、16 件が条件付きで省略、2 件が
既知の期待失敗として扱われた。対象試験は重複なく、最終原稿検証 64、
カーネル・遷移・期・状態・創設・方式 285、敵対検査・統治 126、修復・
隔離再生成 66、効用進化 85、評価例・測定 157、指示文進化 93、論文役
94、文脈・予算・メタ進化・提案・接続処理 154 件に分かれ、合計は
1,124 件である。機械可読な成果物には全試験識別子を収録した。構成要素の遷移表、権限行列、再開
手続きとの対応は \cref{sec:appendix:specification} に示した。

以前記載した完全長のコミット識別子は公開コード保管庫で照合できず、未確定の本作業版
も指していないため、公開成果物という主張を撤回する。この作業木には、
ソースの要約値、依存関係一覧、正確な対象選択条件、全試験識別子、対象
試験と全体試験の記録、および全組合せを展開した機械可読な権限行列を持つ
局所的な内容識別成果物を追加した。未確定の
作業木を現在の Git 識別子だけで表さず、内容全体の要約値で識別する。
ただし、変更不能な公開タグ、DOI 付き保管物、コンテナ、継続的統合記録
ではなく、第三者が取得できる状態ではない。投稿前に指示文一式とともに
永続公開する必要がある。

二件の期待失敗は、命令層から表示層への依存が一件残ることと、評価処理が
共通のモデル接続処理を介さず提供者ライブラリを直接呼ぶことを表す。
RQGM のカーネル検査を迂回する失敗ではないが、合格には数えず全体記録
へ残している。

試験数は実装規模と回帰防止の証拠であり、科学的性能ではない。現環境
には被覆率計測器を導入しておらず、行・分岐被覆率と変異試験得点は不明
である。すべての不変条件に最低一件の試験があることも、この集計だけ
からは証明できない。したがって、総試験数を安全性の代理指標としては
用いない。

創設件数は構成によって異なる。基本の研究構成は指示文 31 件と効用規則
1 件、構成要素 20 件を登録する。論文保管庫を有効にすると、自己選好
検査、論文査読、論文執筆の各指示文と構成要素が一件ずつ加わり、合計は
指示文等 35 件、構成要素 23 件となる。基本構成要素には登録した研究
生成役を含む。従来の「32 件と 19 件」は訂正前の集計である。

過去の探索的実行では、境界取引、効用規則の交代、保存済み評価軸からの
再採点、論文役の後継採用を観測した。しかし、本報告には完全な保存点、
使用モデルの厳密な版、復号設定、乱数種、独立試行数、課題記述、費用
記録を同梱していない。このため、これらは実装経路を見つけるための観察
であり、再現可能な実験結果として採用しない。

\subsection{故障注入と構成上の保証}\label{sec:eval:repro}

固定された正解付き例では、五種類の故障経路を調べた。実行値から再計算
できない指標操作、証拠のない過大主張、封印対象を含む指示文候補、
隔離再生成の禁止入力、退役指示文に依存する記録の候補集合への混入で
ある。各例は、それぞれ数値・証拠検証、候補検査、憲法カーネル、影響
除去検査の期待した経路で拒否された。一つの正常な数値主張例は誤って
拒否されなかった。

さらに、出来事の種類、識別子、取引識別子、内容、直前ハッシュ値の変更、
確定済み出来事の並べ替え、取引所属の不一致を負例として拒否した。T16 が
現期終了と新しい指紋値の次期開始を行うこと、機械可読な権限一覧が疎行列
の全項目と一致することも検査した。これらは決定論的な回帰試験であり、
無作為な出来事生成や形式的モデル検査ではない。

この結果は、五故障すべてへの検出率 5/5 と正常例 1/1 の通過を、その
固定例の範囲で示すだけである。標本が小さく、故障ごとの難易度も一つ
しかないため、誤検出率や見逃し率の推定値として報告しない。先行研究の
捏造、裁定経路への妨害、責任対象の誤帰属、意味を保った隔離再生成の
汚染は、正解付き検出実験が未完である。

構成上の保証にも境界がある。台帳遷移器以外の書換え、表にない遷移、
候補の権限拡大、未裁定攻撃による減点、未確定取引の再生は固定処理が
拒否する。一方、これは保存済み入力に対するカーネル決定性であり、
モデル出力の再生成性ではない。版二の監査連鎖は、出来事の意味を決める
全項目を結び、試験した途中改変を検出するが、外部固定なしに
保存点全体を古い接頭辞へ戻す末尾切り詰めは検出しない。T16 は現期を
強制終了し、同じ取引で新しい指紋値の期を開始する。

\subsection{限界}\label{sec:eval:limits}

採択水準の実証には、事前登録した課題群、使用モデルと提供者側の版、
温度、標本化法、最大出力量、道具版、乱数種、独立試行数、期数、ノード
数、候補数、攻撃・動議・裁定・退役件数、トークン・時間・金額を保存
する必要がある。各比率には分母、試行間分散、信頼区間を添える。

RQGM 元論文との対応を優先した論文比較用の五条件を実装した。P0 は執筆役
だけを進化させて査読役を固定し、P1 は査読役の交代だけを加え、P2 は
共進化するが旧査読役の効用を残し、P3 はその効用を選択対象外にする
完全な RQGM 条件、P4 は P3 に立憲カーネルの強制を加える。全条件は同じ
論文保管庫、計算予算、乱数種、モデル設定を使う評価用設定であり、新しい
運用方式ではない。主比較は P0、P3、P4、機構の切り分けは P1、P2 で行う。
執筆と独立した固定評価には別々に固定したモデルを割り当てられ、実験では
Claude Code を執筆、Codex を評価基準に基づく事後査読へ用いる。最終稿は
共進化する査読役とは別の固定パネルが採点する。ただし、複数乱数種の
実モデル実行は本稿に含まれず、これは実装済みの評価手順であって結果では
ない。また、機構は元論文に合わせたが、出荷時の予算は8期、各期の展開
上限12件であり、元論文の12,288回評価という計算規模
\cite{iacob2026red}を再現してはいない。投稿用実験では共通の評価呼出し
予算を事前登録し、条件ごとの実績を報告する必要がある。

追加の構成要素別評価では、過去事例再試験、責任帰属、隔離再生成、可変
効用規則、基準資料、固定原稿検証器を一つずつ除く必要がある。責任帰属と
影響修復には人工的な正解付き因果グラフを用い、対象精度、到達範囲の
適合率・再現率、過剰除去を測る必要がある。

現時点で正当化できる結論は限定的である。固定した五故障例では期待した
検査が発火し、出来事完全性の負例を拒否し、1,124 件の関連試験は合格した。
それ以上の研究品質向上、
一般的な検出性能、長期活動性、費用対効果は未実証である。
